% Appendix E: 数学公式集
% 既存の本文で \appendix を宣言し、付録 A--D に続けて \input する。
% 必要なパッケージ: amsmath, amssymb, booktabs, longtable, titlesec, needspace
% Compact layout settings (scoped to this appendix only).
% Designed for the book's A4 layout and an optional A5 handout.
\begingroup
% 各数式の先頭を左揃えにするため、以下の数式には flalign* を使用する。
% & は左端のアンカーであり、等号の位置を揃える目的には使わない。
\titlespacing*{\chapter}{0pt}{0pt}{6pt}
\titlespacing*{\section}{0pt}{0.45em}{0.12em}
\titlespacing*{\subsection}{0pt}{0.32em}{0.08em}
\small
\linespread{1.0}\selectfont
\setlength{\parskip}{0pt}
\setlength{\abovedisplayskip}{3pt plus 1pt minus 1pt}
\setlength{\belowdisplayskip}{3pt plus 1pt minus 1pt}
\setlength{\abovedisplayshortskip}{2pt plus 1pt minus 1pt}
\setlength{\belowdisplayshortskip}{2pt plus 1pt minus 1pt}
\setlength{\jot}{1.5pt}
\setlength{\tabcolsep}{3.5pt}
\renewcommand{\arraystretch}{0.97}
\setlength{\LTpre}{3pt}
\setlength{\LTpost}{3pt}
\allowdisplaybreaks[2]
\chapter{数学公式集}\label{app:math-formulas}

この付録では、本文および付録の計算でよく使う数学公式をまとめる。
特に断らない限り、$\log x$ は自然対数、角度はラジアンで表す。
$u,v$ は微分可能な関数、$f$ は必要な階数まで微分可能な関数とする。
ベクトルの内積・外積にはそれぞれ $\cdot,\times$ を用いる。

\section{基本的な記号、指数・対数}\label{sec:app-math-basic}
\subsection{階乗と二項定理}
\begin{flalign*}
 & n!= n(n-1)\cdots 2\cdot 1, \qquad 0!=1 && \\
 & \binom{n}{k}= \frac{n!}{k!(n-k)!} \qquad (0\le k\le n) && \\
 & (a+b)^n= \sum_{k=0}^{n}\binom{n}{k}a^{n-k}b^k \qquad (n=0,1,2,\ldots) &&
\end{flalign*}

\subsection{指数と対数}
\begin{flalign*}
 & a^{x+y}=a^xa^y && \\
 & a^{-x}=\frac{1}{a^x} && \\
 & \log(xy)=\log x+\log y && \\
 & \log\frac{x}{y}=\log x-\log y && \\
 & \log(x^p)=p\log x && \\
 & \log_a x=\frac{\log x}{\log a} && \\
 & e^{\log x}=x\quad(x>0) && \\
 & \log(e^x)=x &&
\end{flalign*}
ここで、対数の引数は正、$a>0$ とし、対数の底として $a\ne1$ を仮定する。
また、$\log_{10} x=(\log x)/(\log 10)$ である。

\section{三角関数と双曲線関数}\label{sec:app-math-trig}
\subsection{三角関数の基本公式}
\begin{flalign*}
 & \sin^2x+\cos^2x=1 && \\
 & 1+\tan^2x=\frac{1}{\cos^2x} && \\
 & \sin(x\pm y)=\sin x\cos y\pm\cos x\sin y && \\
 & \cos(x\pm y)=\cos x\cos y\mp\sin x\sin y && \\
 & \tan(x\pm y)=\frac{\tan x\pm\tan y}{1\mp\tan x\tan y} &&
\end{flalign*}
\begin{flalign*}
 & \sin x+\sin y=2\sin\frac{x+y}{2}\cos\frac{x-y}{2} && \\
 & \sin x-\sin y=2\cos\frac{x+y}{2}\sin\frac{x-y}{2} && \\
 & \cos x+\cos y=2\cos\frac{x+y}{2}\cos\frac{x-y}{2} && \\
 & \cos x-\cos y=-2\sin\frac{x+y}{2}\sin\frac{x-y}{2} &&
\end{flalign*}
\begin{flalign*}
 & 2\sin x\sin y=\cos(x-y)-\cos(x+y) && \\
 & 2\cos x\cos y=\cos(x-y)+\cos(x+y) && \\
 & 2\sin x\cos y=\sin(x+y)+\sin(x-y) && \\
 & \sin 2x=2\sin x\cos x && \\
 & \cos 2x=\cos^2x-\sin^2x=1-2\sin^2x=2\cos^2x-1 && \\
 & \sin 3x=3\sin x-4\sin^3x && \\
 & \cos 3x=4\cos^3x-3\cos x &&
\end{flalign*}
角度の換算は $180^\circ=\pi\,\mathrm{rad}$ である。

\subsection{双曲線関数}
\begin{flalign*}
 & \sinh x=\frac{e^x-e^{-x}}{2} && \\
 & \cosh x=\frac{e^x+e^{-x}}{2} && \\
 & \tanh x=\frac{\sinh x}{\cosh x} && \\
 & \cosh^2x-\sinh^2x=1 && \\
 & \sinh(x\pm y)=\sinh x\cosh y\pm\cosh x\sinh y && \\
 & \cosh(x\pm y)=\cosh x\cosh y\pm\sinh x\sinh y && \\
 & \sinh 2x=2\sinh x\cosh x && \\
 & \cosh 2x=2\cosh^2x-1=1+2\sinh^2x &&
\end{flalign*}
さらに、逆双曲線関数は\(\operatorname{arsinh}x=\log(x+\sqrt{1+x^2})\)、\(\operatorname{arcosh}x=\log(x+\sqrt{x^2-1})\)（$x\ge1$）と表せる。

\section{微分の公式}\label{sec:app-math-diff}
\subsection{基本法則}
ダッシュは $x$ による微分を表す。
\begin{flalign*}
 & (au+bv)'=au'+bv' && \\
 & (uv)'=u'v+uv' && \\
 & \left(\frac{u}{v}\right)'=\frac{u'v-uv'}{v^2} && \\
 & (u^p)'=pu^{p-1}u' && \\
 & \frac{\mathrm d}{\mathrm dx}f(g(x))=f'(g(x))g'(x) &&
\end{flalign*}
最後の公式は合成関数の微分（連鎖律）である。

\subsection{初等関数の導関数}
\begin{center}\small
\begin{tabular}{@{}ll@{\qquad}ll@{}}
\toprule
$f(x)$ & $f'(x)$ & $f(x)$ & $f'(x)$\\
\midrule
$x^p$ & $px^{p-1}$ & $e^{ax}$ & $ae^{ax}$\\
$a^x$ & $a^x\log a$ & $\log|x|$ & $1/x$\\
$\sin x$ & $\cos x$ & $\cos x$ & $-\sin x$\\
$\tan x$ & $1/\cos^2 x$ & $\arctan x$ & $1/(1+x^2)$\\
$\arcsin x$ & $1/\sqrt{1-x^2}$ & $\arccos x$ & $-1/\sqrt{1-x^2}$\\
$\sinh x$ & $\cosh x$ & $\cosh x$ & $\sinh x$\\
$\tanh x$ & $1/\cosh^2 x$ & $\sqrt{1+x^2}$ & $x/\sqrt{1+x^2}$\\
\bottomrule
\end{tabular}
\end{center}
適用範囲では各関数の引数・分母・平方根の条件を満たすものとする。

\section{積分の公式}\label{sec:app-math-int}
\subsection{基本法則と不定積分}
\begin{flalign*}
 & \int (au+bv)\,\mathrm dx=a\int u\,\mathrm dx+b\int v\,\mathrm dx && \\
 & \int uv'\,\mathrm dx=uv-\int u'v\,\mathrm dx && \\
 & \int f(g(x))g'(x)\,\mathrm dx=\int f(t)\,\mathrm dt\quad(t=g(x)) &&
\end{flalign*}
\begin{flalign*}
 & \int x^p\,\mathrm dx=\frac{x^{p+1}}{p+1}+C\quad(p\ne-1) && \\
 & \int\frac{\mathrm dx}{x}=\log|x|+C && \\
 & \int e^{ax}\,\mathrm dx=\frac{e^{ax}}{a}+C\quad(a\ne0) && \\
 & \int\sin x\,\mathrm dx=-\cos x+C && \\
 & \int\cos x\,\mathrm dx=\sin x+C && \\
 & \int x\sin x\,\mathrm dx=\sin x-x\cos x+C && \\
 & \int x\cos x\,\mathrm dx=\cos x+x\sin x+C && \\
 & \int\frac{\mathrm dx}{x^2+a^2}=\frac{1}{a}\arctan\frac{x}{a}+C\quad(a>0) && \\
 & \int\frac{\mathrm dx}{\sqrt{a^2-x^2}}=\arcsin\frac{x}{a}+C\quad(a>0) && \\
 & \int\frac{\mathrm dx}{\sqrt{x^2+a^2}}=\log\left|x+\sqrt{x^2+a^2}\right|+C &&
\end{flalign*}
指数関数と三角関数の積については、$a^2+b^2>0$ として
\begin{flalign*}
 & \int e^{ax}\sin bx\,\mathrm dx=\frac{e^{ax}}{a^2+b^2}(a\sin bx-b\cos bx)+C && \\
 & \int e^{ax}\cos bx\,\mathrm dx=\frac{e^{ax}}{a^2+b^2}(a\cos bx+b\sin bx)+C &&
\end{flalign*}

\subsection{代表的な定積分とガウス積分}
$a>0$ のとき、
\begin{flalign*}
 & \int_0^\infty e^{-ax}\,\mathrm dx=\frac{1}{a} && \\
 & \int_0^\infty\frac{a\,\mathrm dx}{x^2+a^2}=\frac{\pi}{2} && \\
 & \int_0^\infty e^{-ax^2}\,\mathrm dx=\frac12\sqrt{\frac{\pi}{a}} && \\
 & \int_{-\infty}^{\infty}e^{-ax^2}\,\mathrm dx=\sqrt{\frac{\pi}{a}} && \\
 & \int_0^\infty x^2e^{-ax^2}\,\mathrm dx=\frac14\sqrt{\frac{\pi}{a^3}} && \\
 & \int_0^\infty e^{-ax^2}\cos(bx)\,\mathrm dx=\frac12\sqrt{\frac{\pi}{a}}e^{-b^2/(4a)} &&
\end{flalign*}

\section{テイラー展開}\label{sec:app-math-taylor}
\subsection{一般式と近似誤差}
$h=x-a$ とすると、十分滑らかな $f$ に対して
\begin{flalign*}
 & f(a+h)=f(a)+hf'(a)+\frac{h^2}{2!}f''(a)+\cdots+\frac{h^n}{n!}f^{(n)}(a)+R_{n+1} && \\
 & R_{n+1}=\frac{h^{n+1}}{(n+1)!}f^{(n+1)}(\xi), \qquad \xi\text{ は }a\text{ と }a+h\text{ の間} &&
\end{flalign*}
$f^{(n+1)}$ が近傍で有界であれば $R_{n+1}=O(h^{n+1})$ である。

\subsection{主なマクローリン展開}
\begin{flalign*}
 & e^x=1+x+\frac{x^2}{2!}+\frac{x^3}{3!}+\cdots && \\
 & \sin x=x-\frac{x^3}{3!}+\frac{x^5}{5!}-\cdots && \\
 & \cos x=1-\frac{x^2}{2!}+\frac{x^4}{4!}-\cdots && \\
 & \sinh x=x+\frac{x^3}{3!}+\frac{x^5}{5!}+\cdots && \\
 & \cosh x=1+\frac{x^2}{2!}+\frac{x^4}{4!}+\cdots && \\
 & \frac{1}{1+x}=1-x+x^2-x^3+\cdots\quad(|x|<1) && \\
 & \log(1+x)=x-\frac{x^2}{2}+\frac{x^3}{3}-\cdots\quad(-1<x\le1) && \\
 & (1+x)^p=1+px+\frac{p(p-1)}{2!}x^2+\cdots\quad(|x|<1) && \\
 & \arctan x=x-\frac{x^3}{3}+\frac{x^5}{5}-\cdots\quad(|x|\le1) && \\
 & \tan x=x+\frac{x^3}{3}+\frac{2x^5}{15}+O(x^7) &&
\end{flalign*}
最後の $\tan x$ の式は $|x|<\pi/2$ で収束するマクローリン展開の一部であり、$O(x^7)$ は $x\to0$ における剰余項を表す。

\section{ベクトル演算と積分定理}\label{sec:app-math-vector}
\subsection{内積・外積と三重積}
\begin{flalign*}
 & \mathbf A\cdot\mathbf B=A_xB_x+A_yB_y+A_zB_z && \\
 & (\mathbf A\times\mathbf B)_x=A_yB_z-A_zB_y && \\
 & (\mathbf A\times\mathbf B)_y=A_zB_x-A_xB_z && \\
 & (\mathbf A\times\mathbf B)_z=A_xB_y-A_yB_x && \\
 & \mathbf A\cdot(\mathbf B\times\mathbf C)=\mathbf B\cdot(\mathbf C\times\mathbf A) && \\
 & \mathbf A\times(\mathbf B\times\mathbf C)=\mathbf B(\mathbf A\cdot\mathbf C)-\mathbf C(\mathbf A\cdot\mathbf B) && \\
 & (\mathbf A\times\mathbf B)\cdot(\mathbf C\times\mathbf D)=(\mathbf A\cdot\mathbf C)(\mathbf B\cdot\mathbf D)-(\mathbf A\cdot\mathbf D)(\mathbf B\cdot\mathbf C) &&
\end{flalign*}

\subsection{勾配、発散、回転とベクトル恒等式}
$\phi,\psi$ をスカラー場、$\mathbf A,\mathbf B$ をベクトル場とする。
\begin{flalign*}
 & \nabla(\phi\psi)=\phi\nabla\psi+\psi\nabla\phi && \\
 & \nabla\cdot(\phi\mathbf A)=\nabla\phi\cdot\mathbf A+\phi\nabla\cdot\mathbf A && \\
 & \nabla\times(\phi\mathbf A)=\nabla\phi\times\mathbf A+\phi\nabla\times\mathbf A && \\
 & \nabla\cdot(\mathbf A\times\mathbf B)=\mathbf B\cdot(\nabla\times\mathbf A)-\mathbf A\cdot(\nabla\times\mathbf B) && \\
 & \nabla\times(\mathbf A\times\mathbf B)=(\mathbf B\cdot\nabla)\mathbf A-\mathbf B(\nabla\cdot\mathbf A)-(\mathbf A\cdot\nabla)\mathbf B+\mathbf A(\nabla\cdot\mathbf B) && \\
 & \nabla(\mathbf A\cdot\mathbf B)=(\mathbf B\cdot\nabla)\mathbf A+(\mathbf A\cdot\nabla)\mathbf B+\mathbf B\times(\nabla\times\mathbf A)+\mathbf A\times(\nabla\times\mathbf B) && \\
 & \nabla\times(\nabla\phi)=\mathbf 0 && \\
 & \nabla\cdot(\nabla\times\mathbf A)=0 && \\
 & \nabla\times(\nabla\times\mathbf A)=\nabla(\nabla\cdot\mathbf A)-\nabla^2\mathbf A && \\
 & \nabla^2(\phi\psi)=\phi\nabla^2\psi+\psi\nabla^2\phi+2\nabla\phi\cdot\nabla\psi &&
\end{flalign*}
ここで $(\mathbf A\cdot\nabla)\mathbf B$ は、直交座標では
$A_x\,\partial\mathbf B/\partial x+A_y\,\partial\mathbf B/\partial y+A_z\,\partial\mathbf B/\partial z$
を表す。

\subsection{積分定理}
$\partial V$ を体積 $V$ の外向き法線を持つ境界、$\partial S$ を面 $S$ の向きと整合する境界曲線とする。
\begin{flalign*}
 & \text{Gauss の定理:}\quad
 \int_V\nabla\cdot\mathbf A\,\mathrm dV=\oint_{\partial V}\mathbf A\cdot\mathrm d\mathbf S && \\
 & \text{Stokes の定理:}\quad
 \int_S(\nabla\times\mathbf A)\cdot\mathrm d\mathbf S=\oint_{\partial S}\mathbf A\cdot\mathrm d\boldsymbol\ell && \\
 & \text{Green の定理:}\quad
 \oint_{\partial D}(P\,\mathrm dx+Q\,\mathrm dy)=\iint_D\left(\frac{\partial Q}{\partial x}-\frac{\partial P}{\partial y}\right)\mathrm dx\,\mathrm dy &&
\end{flalign*}
Green の定理では境界を正の向き（反時計回り）に回る。

\section{座標系における微分演算子}\label{sec:app-math-coordinates}
以下では、球座標の $\theta$ は $+z$ 軸から測る極角、$\phi$ は $xy$ 平面内の方位角である。

\subsection{直交座標 $(x,y,z)$}
\begin{flalign*}
 & \nabla f=\mathbf e_x\frac{\partial f}{\partial x}+\mathbf e_y\frac{\partial f}{\partial y}+\mathbf e_z\frac{\partial f}{\partial z} && \\
 & \nabla\cdot\mathbf A=\frac{\partial A_x}{\partial x}+\frac{\partial A_y}{\partial y}+\frac{\partial A_z}{\partial z} && \\
 & \nabla^2 f=\frac{\partial^2f}{\partial x^2}+\frac{\partial^2f}{\partial y^2}+\frac{\partial^2f}{\partial z^2} &&
\end{flalign*}
\begin{flalign*}
 & \nabla\times\mathbf A=\left(\frac{\partial A_z}{\partial y}-\frac{\partial A_y}{\partial z}\right)\mathbf e_x
 +\left(\frac{\partial A_x}{\partial z}-\frac{\partial A_z}{\partial x}\right)\mathbf e_y +\left(\frac{\partial A_y}{\partial x}-\frac{\partial A_x}{\partial y}\right)\mathbf e_z &&
\end{flalign*}

\subsection{2次元極座標 $(R,\phi)$}
\begin{flalign*}
 & x=R\cos\phi,\quad y=R\sin\phi,\quad \mathrm dS=R\,\mathrm dR\,\mathrm d\phi && \\
 & \mathbf e_R=\cos\phi\,\mathbf e_x+\sin\phi\,\mathbf e_y && \\
 & \mathbf e_\phi=-\sin\phi\,\mathbf e_x+\cos\phi\,\mathbf e_y && \\
 & \nabla f=\mathbf e_R\frac{\partial f}{\partial R}+\mathbf e_\phi\frac{1}{R}\frac{\partial f}{\partial\phi} && \\
 & \nabla\cdot\mathbf A=\frac1R\frac{\partial(RA_R)}{\partial R}+\frac1R\frac{\partial A_\phi}{\partial\phi} && \\
 & \nabla^2f=\frac1R\frac{\partial}{\partial R}\left(R\frac{\partial f}{\partial R}\right)+\frac1{R^2}\frac{\partial^2 f}{\partial\phi^2} && \\
 & (\nabla\times\mathbf A)_z=\frac1R\frac{\partial(RA_\phi)}{\partial R}-\frac1R\frac{\partial A_R}{\partial\phi} &&
\end{flalign*}
位置ベクトル $\mathbf r=R\mathbf e_R$ の時間微分は
\begin{flalign*}
 & \dot{\mathbf r}=\dot R\mathbf e_R+R\dot\phi\mathbf e_\phi && \\
 & \ddot{\mathbf r}=(\ddot R-R\dot\phi^2)\mathbf e_R+(2\dot R\dot\phi+R\ddot\phi)\mathbf e_\phi &&
\end{flalign*}

\subsection{円柱座標 $(R,\phi,z)$}
\begin{flalign*}
 & \mathrm dV=R\,\mathrm dR\,\mathrm d\phi\,\mathrm dz && \\
 & \nabla f=\mathbf e_R\frac{\partial f}{\partial R}+\mathbf e_\phi\frac1R\frac{\partial f}{\partial\phi}+\mathbf e_z\frac{\partial f}{\partial z} && \\
 & \nabla\cdot\mathbf A=\frac1R\frac{\partial(RA_R)}{\partial R}+\frac1R\frac{\partial A_\phi}{\partial\phi}+\frac{\partial A_z}{\partial z} && \\
 & \nabla^2 f=\frac1R\frac{\partial}{\partial R}\left(R\frac{\partial f}{\partial R}\right)+\frac1{R^2}\frac{\partial^2f}{\partial\phi^2}+\frac{\partial^2f}{\partial z^2} &&
\end{flalign*}
\begin{flalign*}
 & \nabla\times\mathbf A=\left(\frac1R\frac{\partial A_z}{\partial\phi}-\frac{\partial A_\phi}{\partial z}\right)\mathbf e_R
 +\left(\frac{\partial A_R}{\partial z}-\frac{\partial A_z}{\partial R}\right)\mathbf e_\phi 
 +\left(\frac1R\frac{\partial(RA_\phi)}{\partial R}-\frac1R\frac{\partial A_R}{\partial\phi}\right)\mathbf e_z &&
\end{flalign*}
軸対称場では $\partial/\partial\phi=0$ とできる。

\subsection{球座標 $(r,\theta,\phi)$}
\begin{flalign*}
 & x=r\sin\theta\cos\phi,\quad y=r\sin\theta\sin\phi,\quad z=r\cos\theta && \\
 & \mathrm dV=r^2\sin\theta\,\mathrm dr\,\mathrm d\theta\,\mathrm d\phi && \\
 & \nabla f=\mathbf e_r\frac{\partial f}{\partial r}+\mathbf e_\theta\frac1r\frac{\partial f}{\partial\theta}+\mathbf e_\phi\frac1{r\sin\theta}\frac{\partial f}{\partial\phi} && \\
 & \nabla\cdot\mathbf A=\frac1{r^2}\frac{\partial(r^2A_r)}{\partial r}+\frac1{r\sin\theta}\frac{\partial(\sin\theta A_\theta)}{\partial\theta}+\frac1{r\sin\theta}\frac{\partial A_\phi}{\partial\phi} &&
\end{flalign*}
\begin{flalign*}
 & \nabla^2f=\frac1{r^2}\frac{\partial}{\partial r}\left(r^2\frac{\partial f}{\partial r}\right)
 +\frac1{r^2\sin\theta}\frac{\partial}{\partial\theta}\left(\sin\theta\frac{\partial f}{\partial\theta}\right)
 +\frac1{r^2\sin^2\theta}\frac{\partial^2f}{\partial\phi^2} &&
\end{flalign*}
ベクトル場の回転は球座標で次のように表される。
\begin{flalign*}
 & (\nabla\times\mathbf A)_r=\frac1{r\sin\theta}\left[\frac{\partial(\sin\theta A_\phi)}{\partial\theta}-\frac{\partial A_\theta}{\partial\phi}\right] && \\
 & (\nabla\times\mathbf A)_\theta=\frac1r\left[\frac1{\sin\theta}\frac{\partial A_r}{\partial\phi}-\frac{\partial(rA_\phi)}{\partial r}\right] && \\
 & (\nabla\times\mathbf A)_\phi=\frac1r\left[\frac{\partial(rA_\theta)}{\partial r}-\frac{\partial A_r}{\partial\theta}\right] &&
\end{flalign*}
球対称スカラー場 $f=f(r)$ では特に
\begin{flalign*}
 & \nabla^2f=\frac1{r^2}\frac{\mathrm d}{\mathrm dr}\left(r^2\frac{\mathrm df}{\mathrm dr}\right)=\frac1r\frac{\mathrm d^2(rf)}{\mathrm dr^2}\quad(r>0) &&
\end{flalign*}

\Needspace{11\baselineskip}
\section{フーリエ級数と積分変換}\label{sec:app-math-fourier}
以下の公式は、必要な可積分性・滑らかさ・収束条件の下で用いる。

\subsection{フーリエ級数}
周期 $L$ の実関数 $f(x)$：
\begin{flalign*}
 & f(x)=\frac{a_0}{2}+\sum_{n=1}^{\infty}\left[a_n\cos\left(\frac{2\pi n x}{L}\right)+b_n\sin\left(\frac{2\pi n x}{L}\right)\right] && \\
 & a_n=\frac{2}{L}\int_{-L/2}^{L/2}f(x)\cos\left(\frac{2\pi n x}{L}\right)\mathrm dx\quad(n\ge0) && \\
 & b_n=\frac{2}{L}\int_{-L/2}^{L/2}f(x)\sin\left(\frac{2\pi n x}{L}\right)\mathrm dx\quad(n\ge1) &&
\end{flalign*}
例：$f(x)=x$ の周期的延長（$-L/2<x<L/2$）：
\begin{flalign*}
 & x=\frac{L}{\pi}\sum_{n=1}^{\infty}\frac{(-1)^{n+1}}{n}\sin\left(\frac{2\pi nx}{L}\right) &&
\end{flalign*}
Parsevalの公式（二乗可積分な実関数）：
\begin{flalign*}
 & \frac{1}{L}\int_{-L/2}^{L/2}|f(x)|^2\,\mathrm dx=\frac{a_0^2}{4}+\frac12\sum_{n=1}^{\infty}(a_n^2+b_n^2) &&
\end{flalign*}

\subsection{フーリエ変換}
フーリエ変換・逆変換（本書の規約）：
\begin{flalign*}
 & \hat f(k)=\int_{-\infty}^{\infty}f(x)e^{-ikx}\,\mathrm dx && \\
 & f(x)=\frac{1}{2\pi}\int_{-\infty}^{\infty}\hat f(k)e^{ikx}\,\mathrm dk &&
\end{flalign*}
微分・畳み込み（微分では $|x|\to\infty$ で境界項が消えるとする）：
\begin{flalign*}
 & \widehat{f'}(k)=ik\hat f(k) && \\
 & \widehat{f''}(k)=-k^2\hat f(k) && \\
 & (f*g)(x)=\int_{-\infty}^{\infty}f(x-y)g(y)\,\mathrm dy && \\
 & \widehat{f*g}(k)=\hat f(k)\hat g(k) &&
\end{flalign*}
Parsevalの公式（二乗可積分な関数）：
\begin{flalign*}
 & \int_{-\infty}^{\infty}|f(x)|^2\,\mathrm dx=\frac{1}{2\pi}\int_{-\infty}^{\infty}|\hat f(k)|^2\,\mathrm dk &&
\end{flalign*}
代表的な変換対（$a>0$、$x_0,k_0$ は実定数）：
\begin{center}\small
\begin{tabular}{@{}ll@{\qquad}ll@{}}
\toprule
$f(x)$ & $\hat f(k)$ & $f(x)$ & $\hat f(k)$\\
\midrule
$e^{-ax^2}$ & $\sqrt{\pi/a}\,e^{-k^2/(4a)}$ & $\delta(x)$ & $1$\\
$e^{-a|x|}$ & $2a/(a^2+k^2)$ & $\delta(x-x_0)$ & $e^{-ikx_0}$\\
$1/(x^2+a^2)$ & $(\pi/a)\,e^{-a|k|}$ & $1$ & $2\pi\delta(k)$\\
 & & $e^{ik_0x}$ & $2\pi\delta(k-k_0)$\\
\bottomrule
\end{tabular}
\end{center}
Diracのデルタ関数の基本公式（$a\ne0$）：
\begin{flalign*}
 & \int_{-\infty}^{\infty}\delta(x-x_0)\,\mathrm dx=1 && \\
 & \int_{-\infty}^{\infty}g(x)\delta(x-x_0)\,\mathrm dx=g(x_0) && \\
 & \delta(ax)=\frac{1}{|a|}\delta(x) && \\
 & \delta(x-x_0)=\frac{1}{2\pi}\int_{-\infty}^{\infty}e^{ik(x-x_0)}\,\mathrm dk &&
\end{flalign*}
\par\begingroup\footnotesize
\noindent\textit{注：}抽出公式では$g$が$x_0$で連続とする。デルタ関数とそのフーリエ積分表示、および変換対の表の右半分は、超関数（分布）の意味で扱う。\par
\endgroup


\subsection{ラプラス変換}
ラプラス変換（$s$ は積分が収束する複素数）：
\begin{flalign*}
 & F(s)=\mathcal L\{f(t)\}(s)=\int_0^\infty f(t)e^{-st}\,\mathrm dt &&
\end{flalign*}
時間微分・畳み込み（$G(s)=\mathcal L\{g(t)\}(s)$）：
\begin{flalign*}
 & \mathcal L\{f'(t)\}(s)=sF(s)-f(0^+) && \\
 & \mathcal L\{f''(t)\}(s)=s^2F(s)-sf(0^+)-f'(0^+) && \\
 & (f*g)(t)=\int_0^t f(\tau)g(t-\tau)\,\mathrm d\tau && \\
 & \mathcal L\{f*g\}(s)=F(s)G(s) &&
\end{flalign*}
代表的な変換対（$a>0$、$\omega\in\mathbb R$、$n=0,1,2,\ldots$、$\operatorname{Re}s>0$）：
\begin{center}\small
\begin{tabular}{@{}ll@{\qquad}ll@{}}
\toprule
$f(t)$ & $F(s)$ & $f(t)$ & $F(s)$\\
\midrule
$1$ & $1/s$ & $\sin(\omega t)$ & $\omega/(s^2+\omega^2)$\\
$t^n$ & $n!/s^{n+1}$ & $\cos(\omega t)$ & $s/(s^2+\omega^2)$\\
$e^{-at}$ & $1/(s+a)$ & & \\
\bottomrule
\end{tabular}
\end{center}
\par\begingroup\footnotesize
\noindent\textit{注：}$f(0^+)$ は右側極限。畳み込みの公式では各変換の収束を仮定する。\par
\endgroup

\subsection{ハンケル変換}
0次ハンケル変換・逆変換（$J_0$ は0次の第1種Bessel関数）：
\begin{flalign*}
 & \tilde f(k)=\int_0^\infty f(R)J_0(kR)R\,\mathrm dR && \\
 & f(R)=\int_0^\infty\tilde f(k)J_0(kR)k\,\mathrm dk &&
\end{flalign*}
2次元フーリエ変換との関係（軸対称関数、$k=|\mathbf k|$）：
\begin{flalign*}
 & \hat f_{\mathrm{2D}}(k)=2\pi\tilde f(k) &&
\end{flalign*}
代表的な変換対（$a,R_d>0$）：
\begin{center}\small
\begin{tabular}{@{}ll@{}}
\toprule
$f(R)$ & $\tilde f(k)$\\
\midrule
$e^{-aR^2}$ & $e^{-k^2/(4a)}/(2a)$\\
$\Sigma_0 e^{-R/R_d}$ & $\Sigma_0R_d^2/[1+(kR_d)^2]^{3/2}$\\
\bottomrule
\end{tabular}
\end{center}
\par\begingroup\footnotesize
\noindent\textit{注：}後者は指数円盤の面密度$\Sigma(R)=\Sigma_0e^{-R/R_d}$に対応する。\par
\endgroup

\Needspace{11\baselineskip}
\section{行列の基本演算}\label{sec:app-math-matrix}
$A,B$を積のサイズが適合する行列、$\alpha$を実数、$I$を単位行列とする。
\begin{flalign*}
 & (A+B)_{ij}=A_{ij}+B_{ij},\qquad (\alpha A)_{ij}=\alpha A_{ij} && \\
 & (AB)_{ij}=\sum_{k=1}^{n}A_{ik}B_{kj}\qquad(A:m\times n,\ B:n\times p) && \\
 & (A^{\mathsf T})_{ij}=A_{ji},\qquad (AB)^{\mathsf T}=B^{\mathsf T}A^{\mathsf T} && \\
 & \det(AB)=\det A\det B && \\
 & \operatorname{tr}(AB)=\operatorname{tr}(BA),\qquad \operatorname{tr}A=\sum_i A_{ii} && \\
 & AA^{-1}=A^{-1}A=I,\qquad (AB)^{-1}=B^{-1}A^{-1} &&
\end{flalign*}
和は同じ大きさの行列で定義する。$A^{\mathsf T}$は転置行列であり、一般に$AB\ne BA$である。
行列式では$A,B$を同じ大きさの正方行列、逆行列ではさらに正則とする。$\det A\ne0$は$A$が正則であることと同値である。

\endgroup
